\chapter{Preface}

This book was made for Matt. If others find it interesting or useful, so much
the better. But it was not written for posterity. Works like this one, like the
software they describe, are disposable: made with care, for a purpose, and
without the assumption that they will last (Chapter~\ref{ch:transport}).

The book is based on two kinds of source.

The first is the record itself: the code, design documents, specifications and
commit histories of Silica, of Project BEES, and of the other repositories that
appear in these pages---SwErl, BEAM\_SEMP, the research behind the 2024 paper
on unsafe impedance, and the early Organism draft. The book describes those
projects as they stood in November 2026. Silica changes quickly, and some of
what is described here as planned may since have been built, while some of
what is described as built may since have been replaced.

The second is the reminiscences of the person who envisioned Silica and
oversaw its creation, gathered in a series of interviews. Where the record
shows what was done, the interviews supply why: the doubts, convictions and
turning points that a repository does not keep.

Where the two sources speak to the same event, they have been checked against
each other. Where they differ, the difference is noted. The book's concern
throughout is not how to write Silica code but the ideas behind it, and the
opportunities its design opens up.
